% Auto-generated by analysis/make_table_cec_ablation.py (results/entrada284_contrasts.csv, entrada285_contrasts.csv)
\begin{table}[t]
\caption{Components of the IG, with the action drawn at random among five destruction sizes (no learner; P$||\Cmax$, $12$ classes $\times$ $100$ instances, five seeds). G: greedy reconstruction; R: the removed jobs reinserted on random machines; $0$, $1$, $2$: no descent, the IG's short descent (at most two improving moves), a descent to a local optimum; R2 is the classical ILS and R2* the same with the best fixed kick size of each class, chosen on the test instances. Gain of the first-named variant in points of the makespan (mean over classes) and number of classes (of $12$) in which it wins beyond the margin. Mean gap to the bound: G0 $0.27\%$, R0 $2.07\%$, G1 $0.20\%$, R1 $1.10\%$, G2 $0.19\%$, R2 $0.94\%$, R2* $0.73\%$.}
\label{tab:ablation}
\centering
{\scriptsize\setlength{\tabcolsep}{2.6pt}
\begin{tabular}{@{}lrrrr@{}}
\toprule
 & \multicolumn{2}{c}{Equal iterations} & \multicolumn{2}{c}{Equal CPU time} \\\cmidrule(lr){2-3}\cmidrule(l){4-5}
 & gain & classes & gain & classes \\
\midrule
G0 vs R0: greedy vs random reinsertion & +1.73 & 12 & +1.70 & 12 \\
G1 vs R1: the same, short descent & +0.88 & 10 & +0.92 & 10 \\
G2 vs R2: the same, full descent & +0.73 & 9 & +0.85 & 10 \\
R2 vs R1: full vs short descent, random & +0.16 & 6 & -- & -- \\
G0 vs R2: plain IG vs classical ILS & +0.65 & 8 & +0.79 & 9 \\
G0 vs R2*: the same, best kick size & +0.45 & 7 & +0.57 & 7 \\
\bottomrule
\end{tabular}}
\end{table}
